\subsection{From the Perceptron to Deep Networks}
Before introducing the architecture of a feedforward neural network, it is useful to recall the
building block from which it originated: the artificial neuron, also known as the \textbf{perceptron}.

\subsubsection{Electronic Brain}
In the first perceptron model, the inputs form a vector of \textbf{binary features} that are simply summed together and passed through a \textbf{threshold function},
without any \textbf{weights} involved.
\begin{figure}[H]
    \centering
    \scalebox{0.8}{
        \begin{tikzpicture}[
  node distance=0.8cm,
  input/.style={circle, draw=blue!80!black, fill=blue!10, minimum size=20pt, font=\small},
  neuron/.style={circle, draw=blue!80!black, fill=blue!20, minimum size=25pt, font=\small\bfseries},
  output/.style={circle, draw=blue!80!black, fill=blue!30, minimum size=30pt, font=\small\bfseries},
  weight/.style={font=\normalfont, color=blue!70!black}
]

  % Input layer nodes
  \node[input] (x1) {$x_1$};
  \node[input, below=of x1] (x2) {$x_2$};
  \node[below=0.2cm of x2] (dots) {$\vdots$};
  \node[input, below=0.2 cm of dots] (xn) {$x_n$};

  % Summation node
  \node[neuron, right=1.5cm of x2] (sum) {$\Sigma$};

  % Activation function node
  \node[neuron, right=1cm of sum] (activation) {$\sigma$};

  % Output node
  \node[output, right=1cm of activation] (output) {$\hat{y}$};

  % Edges from inputs to summation with weight labels
  \draw[->, thick] (x1) -- (sum);
  \draw[->, thick] (x2) -- (sum);
  \draw[->, thick] (xn) -- (sum);

  % Edge from summation to activation
  \draw[->, thick] (sum) -- (activation);

  % Edge from activation to output
  \draw[->, thick] (activation) -- (output);

\end{tikzpicture}

    }
    \caption{Architecture of the electronic brain model. Inputs $x_i$ are summed together and passed through an activation function $\sigma$ to produce the output $\hat{y}$.}
    \label{fig:electronic_brain}
\end{figure}
It can be shown that this model is capable of classifying linearly separable data, such as the AND, OR, and NOT functions, but it fails to classify non-linearly separable data, such as the XOR function.

\subsubsection{The Perceptron}
The \textbf{Perceptron} is an extension of the electronic brain model.
The key difference is the introduction of \textbf{weights} for each input, which allows the model to learn from data and adjust its parameters accordingly.
\begin{figure}[H]
    \centering
    \scalebox{0.8}{
        \begin{tikzpicture}[
  node distance=0.8cm,
  input/.style={circle, draw=neuronedge, fill=neuronfillin, minimum size=20pt, font=\small},
  neuron/.style={circle, draw=neuronedge, fill=neuronfillhid, minimum size=25pt, font=\small\bfseries},
  output/.style={circle, draw=neuronedge, fill=neuronfillout, minimum size=30pt, font=\small\bfseries},
  bias/.style={circle, draw=biasedge, fill=biasfill, minimum size=20pt, font=\small},
  weight/.style={font=\normalsize, color=neuronedge}
]

  % Input layer nodes
  \node[input] (x1) {$x_1$};
  \node[input, below=of x1] (x2) {$x_2$};
  \node[below=0.2cm of x2] (dots) {$\vdots$};
  \node[input, below=0.2 cm of dots] (xn) {$x_n$};

  % Summation node
  \node[neuron, right=1.5cm of x2] (sum) {$\Sigma$};

  % Bias node (above summation)
  \node[bias, above=1cm of sum] (bias) {$1$};

  % Activation function node
  \node[neuron, right=1cm of sum] (activation) {$\sigma$};

  % Output node
  \node[output, right=1cm of activation] (output) {$\hat{y}$};

  % Edges from inputs to summation with weight labels
  \draw[->, thick] (x1) -- (sum) node[midway, above, weight] {$\theta_1$};
  \draw[->, thick] (x2) -- (sum) node[midway, above, weight] {$\theta_2$};
  \draw[->, thick] (xn) -- (sum) node[midway, below, weight] {$\theta_n$};

  % Edge from bias to summation
  \draw[->, thick, biasedge] (bias) -- (sum) node[midway, right, font=\normalsize, color=biasedge] {$\theta_0$};

  % Edge from summation to activation
  \draw[->, thick] (sum) -- (activation);

  % Edge from activation to output
  \draw[->, thick] (activation) -- (output);

\end{tikzpicture}

    }
    \caption{Architecture of a single layer perceptron. Inputs $x_i$ are weighted by $w_i$, summed together, and passed through an activation function $\sigma$ to produce the output $y$.}
    \label{fig:single_perceptron}
\end{figure}
The prediction $\hat{y}$ of the perceptron is given by the following formula:
\begin{equation}
    \hat{y} = \sigma\left(\sum_{i=1}^{n} \theta_i x_i + \theta_0\right)
\end{equation}
where $\theta_i$ are the weights associated with each input $x_i$, and $\sigma$ is the activation function.
We can also express the prediction in vector form as follows:
\begin{equation}
    \hat{y} = \sigma\left(\vct{\theta}^T \vct{x} + \theta_0\right)
\end{equation}

The most important intuition is that the weights $\theta_i$ are not fixed, but rather learned from data.

The bias term is introduced to allow the model to shift the decision boundary (hyperplane) away from the origin.

\termdef{Perceptron}{
    An artificial neuron, i.e., a perceptron is a non-linear parametrized function
    with a restricted output range.
}
Today, the activation function $\sigma$ is typically a non-linear function, such as the sigmoid, ReLU, or tanh function,
which allows the model to learn complex patterns in the data.

\paragraph{Perceptron Learning Algorithm}
\label{perceptron-training}
The weights of the perceptron are learned from data using the \textbf{Perceptron Training Rule}.
Given a set of training examples $\mathcal{T} = \{(\vct{x}_1, y_1), (\vct{x}_2, y_2), \ldots, (\vct{x}_N, y_N)\}$, $y_i \in \{0, 1\}$
and a learning rate $\eta$, the weights are updated as follows:
\begin{enumerate}
    \item Initialize the weights $\theta_i$ to small random values;
    \item Until convergence, repeat:
    \begin{enumerate}
        \item Select randomly a training example $(\vct{x}_i, y_i)$ from $\mathcal{T}$;
        \item Compute the prediction $\hat{y}_i$ using the current weights:
        \[ \hat{y}_i = \sigma\left(\vct{\theta}^T \vct{x}_i + \theta_0 \right) \]
        \item Update the weights based on the error between the prediction and the true label:
        \[ \vct{\theta}^{(t+1)} \leftarrow \vct{\theta}^{(t)} + \eta (y_i - \hat{y}_i) \vct{x}_i \]
        \item Update the bias:
        \[ \theta_0^{(t+1)} \leftarrow \theta_0^{(t)} + \eta (y_i - \hat{y}_i) \]
    \end{enumerate}
\end{enumerate}
Note that we are assuming that the training data is linearly separable, otherwise we would have to introduce a stopping criterion to prevent the algorithm from running indefinitely.

\subsubsection{MLP}
A single perceptron cannot learn non-linearly separable data, such as the XOR function. This limitation motivated the
\textbf{Multilayer Perceptron} (MLP), which stacks several layers of neurons.
\begin{figure}[H]
    \centering
    \scalebox{0.8}{
        \begin{tikzpicture}[node distance=1cm and 2cm]

  % Input layer
  \foreach \i in {1,...,4} {
    \node[neuron in] (i\i) at (0, -\i*1.2) {$x_{\i}$};
  }

  % Hidden layer 1
  \foreach \j in {1,...,3} {
    \node[neuron hid] (h1\j) at (2.5, -0.6 - \j*1.2) {};
  }

  % Hidden layer 2
  \foreach \j in {1,...,3} {
    \node[neuron hid] (h2\j) at (5, -0.6 - \j*1.2) {};
  }

  % Hidden layer 3
  \foreach \j in {1,...,3} {
    \node[neuron hid] (h3\j) at (7.5, -0.6 - \j*1.2) {};
  }

  % Output layer (centered vertically with the network)
  \foreach \k in {1,2} {
    \node[neuron out] (o\k) at (10, -1.2 - \k*1.2) {$y_{\k}$};
  }

  % ──
  \foreach \i in {1,...,4} {
    \foreach \j in {1,...,3} {
      \draw[conn] (i\i) -- (h1\j);
    }
  }

  % ── Connections: Hidden 1 → Hidden 2 ──
  \foreach \i in {1,...,3} {
    \foreach \j in {1,...,3} {
      \draw[conn] (h1\i) -- (h2\j);
    }
  }

  % ── Connections: Hidden 2 → Hidden 3 ──
  \foreach \i in {1,...,3} {
    \foreach \j in {1,...,3} {
      \draw[conn] (h2\i) -- (h3\j);
    }
  }

  % ── Connections: Hidden 3 → Output ──
  \foreach \i in {1,...,3} {
    \foreach \k in {1,2} {
      \draw[conn] (h3\i) -- (o\k);
    }
  }

  % ── Layer labels ──
  \node[layerlabel, above=0.5cm of i1]  {Input\\Layer};
  \node[layerlabel, above=0.5cm of h21] {Hidden Layers};
  \node[layerlabel, above=0.5cm of o1]  {Output\\Layer};

\end{tikzpicture}

    }
    \caption{Architecture of an MLP with three hidden layers, fully connected in a feed-forward fashion.}
    \label{fig:mlp}
\end{figure}
This structure is very flexible: multiple nodes can be used in the output layer, so it can be applied to multi-class classification problems, and also to regression.

\paragraph{Backpropagation}
\label{introduction-backpropagation}
The perceptron learning algorithm is not suitable for training MLPs, as no expected output is available for the hidden layers, but just for the output layer.

Backpropagation is an efficient algorithm for computing the gradients also for large datasets.
At a high level, the algorithm can be summarized as follows:
\begin{enumerate}
    \item \textbf{Forward propagation}: Given an input $\vct{x}$, compute the output of the network $\hat{y}$ by
    propagating the input through the layers of the network;
    
    \item \textbf{Loss computation}: Compute the loss $\loss(y, f(\vct{x}, \vct{\theta}))$, where $y$ is the true label and $f(\vct{x}, \vct{\theta})$ is the output of the network with parameters $\vct{\theta}$;
    
    \item \textbf{Backward propagation}: the error signal is propagated backwards through the network and it's used to update the weights.
    We are basically computing how much each weight contributed to the error, and we can use this information to update the weights in the direction that reduces the error.
\end{enumerate}

This concept is covered in more detail later in this chapter.

\paragraph{Cybenko's Theorem}
Cybenko's theorem states that a 2-layer feed-forward neural network with linear output
can approximate any continuous function on a compact subset of $\R^n$ with arbitrary accuracy.

In practice, the more hidden layers we have, the more complex functions we can learn. Deep
neural networks need a lot of labeled data to be trained, and cannot be trained easily.
